\section{Hypnagogia and the ``Tetris Effect''}\label{sec:hypnagogia}

Before the brain can consolidate and generalize memories, it must first rehearse and replay them.

The first hints of this process occur during the N1 sleep-onset stage, in a state known as \textit{hypnagogia}. This is the twilight period where our directed, waking thoughts begin to fragment and dissolve into spontaneous, dream-like imagery.

A key phenomenon in this stage is the \textbf{``Tetris Effect''}. If an individual spends hours engaged in a novel, repetitive, high-salience activity (like playing Tetris, skiing, or coding), they will often experience spontaneous, intrusive, and fragmented replays of that activity during sleep onset~\cite{stickgold2000tetris}.
\begin{itemize}
    \item An amnesiac patient, who had no conscious memory of playing Tetris, still reported seeing ``falling blocks'' during hypnagogia.
    \item This demonstrates that this ``replay'' is not a conscious, high-level memory recall, but a fundamental, bottom-up playback of recently over-trained procedural or perceptual circuits.
\end{itemize}
This phenomenon is the first and simplest cognitive echo of NREM's function: the brain has tagged a new, high-priority skill and is beginning the process of ``offline'' rehearsal, even before deep sleep is initiated.
